\section{Application}

\begin{UPSTIexercice}{Résistance équivalente (montage série)}
	\UPSTIetape Écrire un programme \texttt{resistance\_serie.c} qui demande à l'utilisateur de saisir des valeurs de résistances (en ohms), une par une, jusqu'à ce qu'il entre \texttt{0} pour arrêter la saisie. Le programme affichera alors la résistance équivalente du montage en série (la somme des résistances saisies). On ne connaît pas à l'avance le nombre de résistances qui seront saisies.
	Par exemple :
	\begin{lstlisting}[language=bash,style=console]
Entrez une resistance en ohms (0 pour arreter) : 220
Entrez une resistance en ohms (0 pour arreter) : 330
Entrez une resistance en ohms (0 pour arreter) : 100
Entrez une resistance en ohms (0 pour arreter) : 0
Resistance equivalente : 650.00 ohms
\end{lstlisting}
\end{UPSTIexercice}

\begin{UPSTIcorrectionP}{Résistance équivalente (montage série)}
	\begin{lstlisting}[language=c]
#include <stdio.h>

int main(void) {
    double resistance;
    double equivalente = 0;

    printf("Entrez une resistance en ohms (0 pour arreter) : ");
    scanf("%lf", &resistance);

    while (resistance != 0) {
        equivalente += resistance;
        printf("Entrez une resistance en ohms (0 pour arreter) : ");
        scanf("%lf", &resistance);
    }

    printf("Resistance equivalente : %.2f ohms\n", equivalente);

    return 0;
}
	\end{lstlisting}
\end{UPSTIcorrectionP}
